\documentclass[a4paper,12pt]{article}
\usepackage[margin=1.5cm]{geometry}
\usepackage{amsmath, amssymb}
\usepackage{array}
\usepackage{multicol}

\usepackage[spanish, bidi=basic, provide=*]{babel}


\usepackage{amsmath}
\usepackage{amssymb}
\usepackage{booktabs}
\usepackage{xcolor}
\usepackage{tcolorbox}
\usepackage{enumitem}

% Definición de colores temáticos de Halloween
\definecolor{halloweenOrange}{HTML}{D9531E}
\definecolor{halloweenPurple}{HTML}{4A154B}
\definecolor{halloweenDark}{HTML}{1A1A1A}
\definecolor{halloweenBg}{HTML}{FDF8F5}

% Configuración de las cajas de retos
\tcbset{
  colback=halloweenBg,
  colframe=halloweenOrange,
  fonttitle=\bfseries\color{white},
  coltitle=white,
  colbacktitle=halloweenOrange,
  boxrule=1pt,
  arc=3mm,
  left=4mm,
  right=4mm,
  top=3mm,
  bottom=3mm
}

\pagestyle{empty}

\begin{document}

% Encabezado
\begin{center}
    \includegraphics[scale=0.001]{../../img/iconoCalabaza.png}
    {\LARGE \textbf{\color{halloweenPurple}  ESCAPE ROOM: LA MALDICIÓN DEL BRUJO EXPONENTIALIS }} \\[0.2cm]
    {\large \textbf{\color{halloweenOrange} Desafío Matemático de Halloween — 3.º / 4.º ESO}} \\[0.4cm]
    \rule{\textwidth}{1.5pt}
\end{center}

\vspace{0.2cm}

\noindent \textbf{Nombre del equipo / Alumno(a):} \dotfill \\[0.2cm]
\noindent \textbf{Curso y grupo:} \dotfill \hfill \textbf{Fecha:} \dotfill

\vspace{0.4cm}

% Historia de ambientación
\begin{tcolorbox}[colframe=halloweenPurple, colbacktitle=halloweenPurple, title= EL CONTEXTO DE LA MISION]
El laboratorio del instituto ha sido dominado por el Gran Brujo \textbf{Exponentialis}. Para romper el hechizo y escapar antes de que caiga la medianoche, debéis descifrar el \textbf{código secreto de 4 dígitos} resolviendo cada una de las pruebas que encontréis a continuación. 

\textbf{¡Atención!} Cada prueba os revelará una cifra del código final. ¡Buena suerte!
\end{tcolorbox}

\vspace{0.4cm}

% Prueba 1
\begin{tcolorbox}[title= PRUEBA 1: El Elixir de las Potencias (Primer Cifra)]
Simplifica la siguiente expresión utilizando las propiedades de las potencias para dejarla en forma de $2^x$:

\[
\frac{(2^3 \cdot 2^{-1})^2 \cdot 2^5}{2^4 \cdot (2^2)^2}
\]

\textbf{Indicación:} Aplica correctamente las propiedades del producto, cociente y potencia de una potencia. \\
\textbf{Dígito de la clave:} El \textbf{primer dígito} del candado es el valor del exponente final $x$.

\vspace{1.5cm}
\hfill \textbf{DÍGITO 1 = } \framebox[1.5cm]{\vrule height 0.5cm depth 0.2cm width 0pt}
\end{tcolorbox}

\vspace{0.4cm}

% Prueba 2
\begin{tcolorbox}[title= PRUEBA 2: El Polvo de Estrellas (Segunda Cifra)]
Para preparar el antídoto contra el hechizo se deben mezclar dos ingredientes medidos a escala microscópica:
\begin{itemize}[leftmargin=1.5em]
    \item \textbf{Veneno de araña negra:} $4{,}2 \cdot 10^{-5}\text{ g}$
    \item \textbf{Pelo de murciélago pulverizado:} $1{,}8 \cdot 10^{-5}\text{ g}$
\end{itemize}

\textbf{Pasos a seguir:}
\begin{enumerate}
    \item Calcula la masa total de la mezcla sumando ambos ingredientes (expresa el resultado en notación científica).
    \item Multiplica la masa total obtenida por el factor de amplificación mágica $10^7$.
    \item Escribe el número resultante en forma estándar.
\end{enumerate}

\textbf{Dígito de la clave:} El \textbf{segundo dígito} es la cifra de las centenas del número final obtenido.

\vspace{1.2cm}
\hfill \textbf{DÍGITO 2 = } \framebox[1.5cm]{\vrule height 0.5cm depth 0.2cm width 0pt}
\end{tcolorbox}

\newpage

% Prueba 3
\begin{tcolorbox}[title= PRUEBA 3: Las Raíces de la Tumba (Tercera Cifra)]
Resuelve la siguiente operación de radicales extraídos de las catacumbas:

\[
\sqrt{75} + 2\sqrt{12} - \sqrt{27}
\]

\textbf{Pasos a seguir:}
\begin{enumerate}
    \item Descompón en factores primos los radicandos y extrae todos los factores posibles fuera de las raíces.
    \item Simplifica la expresión hasta obtener la forma $k\sqrt{3}$.
\end{enumerate}

\textbf{Dígito de la clave:} El \textbf{tercer dígito} de la clave es el número entero $k$ que multiplica a $\sqrt{3}$.

\vspace{2cm}
\hfill \textbf{DÍGITO 3 = } \framebox[1.5cm]{\vrule height 0.5cm depth 0.2cm width 0pt}
\end{tcolorbox}

\vspace{0.6cm}

% Prueba 4
\begin{tcolorbox}[title= PRUEBA 4: La Clave de los Radicales (Cuarta Cifra)]
Aplica las propiedades de los radicales y las potencias de exponente fraccionario para simplificar al máximo la siguiente expresión:

\[
\sqrt{\sqrt[3]{a^{12}}}
\]

\textbf{Pasos a seguir:}
\begin{enumerate}
    \item Utiliza la propiedad de la «raíz de una raíz» ($\sqrt[n]{\sqrt[m]{x}} = \sqrt[n \cdot m]{x}$) o convierte los radicales a potencias de exponente fraccionario.
    \item Simplifica la expresión hasta dejarla de la forma $a^p$.
\end{enumerate}

\textbf{Dígito de la clave:} El \textbf{cuarto dígito} de la clave es el exponente $p$.

\vspace{2cm}
\hfill \textbf{DÍGITO 4 = } \framebox[1.5cm]{\vrule height 0.5cm depth 0.2cm width 0pt}
\end{tcolorbox}

\vspace{0.8cm}

% Caja del Código Final
\begin{center}
\begin{tcolorbox}[colframe=halloweenDark, colback=white, width=0.85\textwidth, halign=center]
    {\Large \textbf{ CÓDIGO FINAL DE ESCAPE}} \\[0.3cm]
    Introduce aquí la combinación de 4 dígitos para desactivar la trampa: \\[0.2cm]
    \textbf{\Large [ \quad \text{—} \quad ] \quad [ \quad \text{—} \quad ] \quad [ \quad \text{—} \quad ] \quad [ \quad \text{—} \quad ]}
\end{tcolorbox}
\end{center}

\end{document}